\documentclass[11pt]{article}

\usepackage[margin=1.0in]{geometry}
\usepackage{array}
\usepackage{booktabs}
\usepackage{tabularx}
\usepackage{longtable}
\usepackage{enumitem}
\usepackage[table]{xcolor}
\usepackage{microtype}
\usepackage{hyperref}
\usepackage{fancyhdr}
\usepackage{titlesec}
\usepackage{ragged2e}

\hypersetup{
    colorlinks=true,
    linkcolor=blue,
    urlcolor=blue,
    citecolor=blue,
    pdftitle={GT 1000 D03 Syllabus},
    pdfauthor={Andrew J. Steinmetz},
    pdfsubject={GT 1000 First-Year Seminar syllabus}
}

\setlength{\headheight}{14pt}
\setlength{\parindent}{0pt}
\setlength{\parskip}{0.35em}
\setlist[itemize]{leftmargin=2em, itemsep=0.15em, topsep=0.25em}
\setlist[enumerate]{leftmargin=2em, itemsep=0.15em, topsep=0.25em}

\definecolor{GTGold}{HTML}{B3A369}
\definecolor{GTNavy}{HTML}{003057}
\definecolor{GTLightGold}{HTML}{F9F6E5} % Diploma
\definecolor{GTGray}{HTML}{54585A}      % Gray Matter
\definecolor{GTLightGray}{HTML}{D6DBD4} % Pi Mile

\titleformat{\section}
  {\large\bfseries\color{GTNavy}}
  {}{0pt}{}
  [\vspace{0.45em}\color{GTGold}\titlerule]

\titleformat{\subsection}
  {\normalsize\bfseries\color{GTNavy}}
  {}{0pt}{}

\titleformat{\subsubsection}
  {\normalsize\bfseries\color{GTNavy}}
  {}{0pt}{}

\titlespacing*{\section}{0pt}{1.2em}{0.6em}
\titlespacing*{\subsection}{0pt}{0.9em}{0.35em}
\titlespacing*{\subsubsection}{0pt}{0.8em}{0.3em}

\pagestyle{fancy}
\fancyhf{}
\fancyhead[L]{\textcolor{black}{GT 1000}}
\fancyhead[R]{\textcolor{black}{First-Year Seminar, Physics}}
\fancyfoot[L]{\textcolor{black}{Georgia Institute of Technology}}
\fancyfoot[R]{\textcolor{black}{\thepage}}
\renewcommand{\headrulewidth}{0.2pt}
\renewcommand{\footrulewidth}{0.2pt}

\newcommand{\dash}{---}
\newcommand{\tablehead}[1]{\textbf{\textcolor{white}{#1}}}
\newcommand{\navycell}[1]{\cellcolor{GTNavy}\textbf{\textcolor{white}{#1}}}

\begin{document}

\noindent\textcolor{GTGold}{\rule{\textwidth}{4pt}}
\vspace{-0.2em}

\begingroup
\setlength{\fboxsep}{14pt}
\noindent\colorbox{GTLightGold}{%
\parbox{\dimexpr\textwidth-2\fboxsep\relax}{%
    {\LARGE\bfseries\textcolor{GTNavy}{GT 1000: First-Year Seminar, Physics}}\\[0.45em]
    {\normalsize\bfseries Fall 2026 \textbar{} CRN \#94631 \textbar{} 1 credit hour}\\[0.25em]
    {\normalsize Fridays, 12:30--1:20 PM \textbullet{} Swann 325
    \textbullet{} Weekly in-person attendance required}
}}
\endgroup

\vspace{0.1em}
\noindent\textcolor{GTLightGray}{\rule{\textwidth}{0.5pt}}

\section*{Course Description}

This seminar course is designed to help you make a successful transition
to college by becoming better acquainted with the academic and social
environments here at Georgia Tech. Through the course, you will acquire
strategies that promote academic, social, and professional success. This
is a highly interactive course that requires active student participation
and working collaboratively in small groups.

\subsection*{Learning Outcomes}

Upon successful completion of this course, students will:
\begin{enumerate}
    \item Develop self-awareness and self-knowledge, particularly as it
    relates to academic interests and professional goals.

    \item Describe their path to graduation and identify opportunities
    for academic and professional enrichment with the assistance of a
    faculty member or advisor.

    \item Demonstrate the ability to lead and productively participate
    in collaborative projects.

    \item Be aware of what it means to be a physicists and explore opportunities for physicists at the undergraduate, graduate and career levels.
\end{enumerate}

\subsection*{Required Course Materials}

None.

\section*{Instructor Information}

\begin{tabularx}{\textwidth}{>{\bfseries}p{0.23\textwidth} X}
\toprule
Instructor & Andrew J. Steinmetz \\
\midrule
Email & \href{mailto:ajsteinmetz@gatech.edu}{ajsteinmetz@gatech.edu} \\
\midrule
Office & Howey W204 \\
\midrule
Office Hours &
\href{https://outlook.office.com/bookwithme/user/118287044b844adb8cb5cfaadd54546f@gatech.edu?anonymous\&ismsaljsauthenabled\&ep=plink}
{By appointment} \\
\bottomrule
\end{tabularx}

\newpage

\section*{Assignments and Grading}

Grading for the course will be broken down as follows:

\begin{center}
\renewcommand{\arraystretch}{1.15}
\begin{tabularx}{0.98\textwidth}{X p{0.18\textwidth}}
\toprule
\textbf{Course Component} & \textbf{Weight} \\
\midrule
Participation, Short Assignments, and Attendance & 25\% \\
Team Project & 25\% \\
Academic Planning & 25\% \\
Career Readiness & 25\% \\
\midrule
\textbf{Total} & \textbf{100\%} \\
\bottomrule
\end{tabularx}
\renewcommand{\arraystretch}{1.0}
\end{center}

\subsection*{Grading Scale}

Your final grade will be assigned as a letter grade according to the
following scale:

\begin{center}
\renewcommand{\arraystretch}{1.15}
\begin{tabularx}{0.98\textwidth}{X p{0.24\textwidth}}
\toprule
\textbf{Letter Grade} & \textbf{Percentage} \\
\midrule
A & 90--100\% \\
B & 80--89\% \\
C & 70--79\% \\
D & 60--69\% \\
F & $<60\%$ \\
\bottomrule
\end{tabularx}
\renewcommand{\arraystretch}{1.0}
\end{center}

Final numeric grades are rounded to the nearest whole percent. For example, 83.76\% becomes 84\%. This course is not otherwise curved. At Georgia Tech, grades are awarded on a scale of A--F with no plus/minus grades permitted.

\href{http://registrar.gatech.edu/info/grading-system}{Georgia Tech grading system}

\section*{Course Policies}

\subsection*{Attendance Policy}

This class is designed for students to become active participants, so
it is vital that all students attend each class session and are fully
engaged in all class activities. This means that you will need to do
more than just show up at the appointed time and place---you will also
need to come to class prepared to participate.

Please attend every scheduled meeting on time. That said, if you are
exhibiting any symptoms of illness (fever over 100.4~F, coughing,
nausea, etc.), please stay home to rest, seek medical attention if
appropriate, and contact me as soon as possible.

Excessive lateness and absences (three or more per semester) can
negatively impact your final grade outcome. If you have extenuating
circumstances, I will try to work with you to address those challenges.
Please communicate early and often if you are struggling with issues
that may call for accommodations.

To request an excused absence, contact the Dean of Students office for class absence verification in the case of a documented illness, hospitalization, accident, death in the family, family emergency, or lengthy illness.

\href{https://studentlife.gatech.edu/dean-students/class-attendance}{Dean of Students: Class Attendance}

\subsubsection*{Late Work}

\textbf{Late work is not accepted.} An assignment may be excused only through the Dean of Students office. While this is not meant to be a difficult course, all assignments are open and available for the entirety of the semester; therefore, appropriate time-management is expected.

\subsection*{Criteria for Successful Completion and Acceptable Student Conduct}

Creating an engaging and productive class experience is the
responsibility of both the instructor and the students involved. To
ensure you are getting the most out of our class and contributing
positively to our learning environment, please be mindful of the
following:

\begin{enumerate}
    \item Please create a habit of checking your GT e-mail and Canvas
    daily. Important class announcements and information will be posted
    to Canvas, and you want to make sure you are fully equipped for
    success.

    \item Grades will be posted to Canvas throughout the semester. It
    is your responsibility to keep track of your submitted assignments
    and grade progress throughout the semester. If something seems
    missing or off to you, please reach out to me to discuss.

    \item GT 1000 is designed to be informative, helpful, and fun.
    However, the class is graded and does carry one credit towards your
    degree. For that reason, it is important that all assignments you
    submit be high quality. Please double-check your work for spelling,
    capitalization, and grammatical errors. This practice is also good
    for creating a habit of professionalism.

    \item You can expect that, as your GT 1000 instructor, I will come
    to class each day ready to engage you in learning about Georgia Tech
    and learning how to be successful academically and professionally.
    Each session will be structured to encourage your active
    participation, so please come prepared to participate and contribute
    to the class.

    \item To ensure everyone feels safe, heard, and able to learn in
    our class, please demonstrate appropriate classroom behavior at all
    times. This means that you should speak and act in a respectful
    manner towards your classmates, your Team Leader, and your
    instructor.
\end{enumerate}

\href{https://policylibrary.gatech.edu/student-life/student-code-conduct}
{Georgia Tech Student Code of Conduct}

\subsection*{Academic Integrity and Generative AI}

GT 1000 is a class based in your own discovery, exploration, and
reflection of what you want your time at Tech and professional life to
be. Unpacking your thought processes based on information you acquire
through campus resources is a major component of the assignments.
Therefore, the use of AI is strongly discouraged.

If you wish to use AI as part of your brainstorming process, that is
fine; however, per the Georgia Tech Honor Code, no work produced via AI
should be submitted as your own. If you have questions about what
constitutes acceptable use, please see the instructor.

The Georgia Tech Academic Honor Code applies to all work submitted in
this course.

\href{https://policylibrary.gatech.edu/student-life/academic-honor-code}
{Georgia Tech Academic Honor Code}

\href{https://oit.gatech.edu/ai/guidance}{Georgia Tech OIT guidance on AI}

\subsection*{Accommodations for Students with Disabilities}

If there are aspects of the instruction or design of this course that
result in barriers to your inclusion or accurate assessment of
achievement, please notify me as soon as possible so we can resolve the
issue.

Students with disabilities should also contact the Office of Disability
Services (ODS), whose purpose is to collaborate with students, faculty,
and staff to create a campus environment that ensures all students have
an equal opportunity to access the Georgia Tech community. ODS can be
reached at \href{tel:4048942563}{404-894-2563},
\href{mailto:dsinfo@gatech.edu}{dsinfo@gatech.edu}, or through the
Office of Disability Services website.

Please contact me ahead of time to discuss any issues related to
disabilities. I am happy to work with you.

\href{https://disabilityservices.gatech.edu/}
{Georgia Tech Office of Disability Services}

\subsection*{Student-Faculty Expectations}

At Georgia Tech, we believe that it is important to strive for an atmosphere of mutual respect, acknowledgment, and responsibility between faculty members and the student body.

In the end, simple respect for knowledge, hard work, and cordial interactions will help build the environment we seek. Therefore, I encourage you to remain committed to the ideals of Georgia Tech while in this class.

\href{http://www.catalog.gatech.edu/rules/21/}{Georgia Tech student-faculty expectations}

\end{document}